\subsection{Term scoring measures length; theme scoring does not}\label{app:representation}

Figure~\ref{fig:representation} puts the two representations side by side
over the same axes, with each hexagon coloured by the mean length of the
filings in it. Under term scoring the colour runs from dark near the
origin --- a median of about 200 terms --- to near-white at the far end,
where the median is eight. That gradient is the measurement problem: a
filing scores as atypical largely because it is short. Under theme scoring
the same surface is a narrow band of roughly uniform tone.

Two filings are marked in both columns. Under the production scoring they
behave differently from one another, which is the point. \textsc{Amazon.com}
(1997) is at the 98th percentile of lead on terms and the 94th on themes,
but only the 40th and 47th
percentiles of atypicality, so both agree that it was early and neither
mistakes it for strange. Apple's \textsc{iPod} (2001) divides them: the 31st
percentile of lead on terms against the 62nd on themes. The surfaces plotted
here are on the class-year build the figure was generated from, where the
length gradient is starkest; the percentiles just quoted are on the per-filing
build used for every estimate.

\begin{figure}[h]\centering
\includegraphics[width=\textwidth]{results/representation_appendix.png}
\caption{Mean filing length across the two axes, by representation and
class, for filings 1995--2019 scored on per-filing reference windows. The
horizontal axis is surprise against the preceding five years, the vertical
axis surprise against the following five; a filing on the dashed diagonal
is equally far from both, and distance below it is lead. Hexagons holding fewer than 25 filings are blank; the retained
hexagons cover over 99.4\% of filings in every panel. The colour scale is
shared across panels; the axis ranges are not, because term-scored and
theme-scored KL are not on a common scale, so only shape and colour
gradient are comparable across columns. Both
columns are drawn on the filings the per-filing windows score, which
differ between representations by under one percent.}
\label{fig:representation}
\end{figure}

Each representation has the virtue the other lacks: theme scoring is
independent of document length, which term scoring is not, while term
scoring resolves novel word \emph{combinations}, which theme scoring does
not. We therefore use theme scoring for every estimate, because the
length confound contaminates firm-level comparisons, and term scoring for
every worked example, because the theme basis has no coordinate for the
language the examples turn on.

\subsection{How three filings encode}\label{app:exhibit}

Table~\ref{tab:exhibit} traces three filings through term scoring and theme
scoring at $T \in \{50, 200, 500\}$ --- three kinds of novelty, and how each
representation handles them.

\textsc{Amazon.com} (1997, advertising/retail) is \emph{combinatorial}
novelty: every word is ordinary, but the bigrams assemble an invention
(``computerized line,'' ``line search,'' ``search ordering''; class
document frequencies 379, 134, 355). The filing's text reads ``on line''
as two words, the standard spelling in 1997, before \emph{online}
lexicalized into a single token; the analyzer drops the stopword ``on''
and the bigram ``line search'' survives
stopword removal. The novelty was necessarily combinatorial
because the word for it did not yet exist. Term scoring sees these
bigrams; under theme scoring at every resolution all of them are out of
vocabulary and the filing reads as a media retailer that uses computers.

\textsc{Kombucha} (1997, beer/soft drinks) is \emph{lexical} novelty: the
new thing arrives as new words (\emph{kombucha} df 567, \emph{saccharose}
out of vocabulary even at the class level). The word \emph{kombucha} is in
the theme vocabulary but, with no kombucha-like theme, the filing loads on
a generic beverages theme at every $T$ (Table~\ref{tab:exhibit}).

\textsc{Ethereum} (2018, software) is \emph{wave} novelty: emergent
vocabulary already diffusing (\emph{blockchains}, \emph{distributed
computing}). It loads on generic software and services themes at every
resolution, including $T = 200$, which does contain a
cryptocurrency/blockchain theme; there, the eleven occurrences of
\emph{software} dominate the theme assignment.

\begin{table}[h]\centering\small
\caption{Top theme loading for three filings, by representation. Term
scoring resolves the novel content of all three; theme scoring resolves
none of it at any tested resolution.}
\label{tab:exhibit}
\begin{tabular}{lp{3.1cm}p{3.1cm}p{3.1cm}}
\toprule
& \textsc{Amazon.com} & \textsc{Kombucha} & \textsc{Ethereum} \\
& (1997) & (1997) & (2018) \\
\midrule
novel content (term) & line search, search ordering & kombucha, saccharose & blockchains, distributed computing \\
$T=50$ & video, computer, music (0.40) & beverages, fruit, drinks (0.64) & software, online, downloadable (0.33) \\
$T=200$ & entertainment, video, games (0.29) & beverages, drinks, coffee (0.60) & field, services, development (0.37) \\
$T=500$ & computer, software, data (0.22) & beverages, drinks, fruit (0.61) & software, online, computer (0.49) \\
\bottomrule
\end{tabular}
\end{table}

The pattern across the row is the measurement point in miniature: theme
scoring cannot represent combinatorial novelty (Amazon), dilutes lexical
novelty (Kombucha), and majority-votes away wave novelty (Ethereum) even
when a fitting theme exists. That the mark-level outcome results are
nearly identical under both scorings despite this blindness is what
licenses theme scoring as a robustness instrument; that the firm-level
results diverge is what exposes them as artifacts of the lexical
specificity term scoring additionally captures
(Section~\ref{sec:robust_hazard}).

